\documentclass[aspectratio=169,12pt]{beamer}
\usepackage{amsmath}
\usepackage{amssymb}
\usepackage{array}
\usepackage[most]{tcolorbox}
\usepackage{graphicx}
\usepackage[sfdefault,lf]{FiraSans}
\renewcommand{\familydefault}{\sfdefault}
\setlength{\parindent}{0pt}
\everymath{\displaystyle}

\setbeamertemplate{navigation symbols}{}
\setbeamertemplate{footline}{}
\setbeamertemplate{headline}{}
\setbeamertemplate{blocks}[rounded][shadow=false]

\definecolor{mmpageWhite}{HTML}{FCFBF7}
\definecolor{mmtanSoft}{HTML}{F7F5EF}
\definecolor{mmgreenDeep}{HTML}{244B2C}
\definecolor{mmgreenLeaf}{HTML}{5D8A56}
\definecolor{mmgreenSoft}{HTML}{E4EFD9}
\definecolor{mmtext}{HTML}{163221}
\definecolor{mmwarn}{HTML}{8F2F36}
\definecolor{mmwarnSoft}{HTML}{F8E8EA}

\setbeamercolor{background canvas}{bg=mmpageWhite}
\setbeamercolor{normal text}{fg=mmtext,bg=mmpageWhite}
\setbeamercolor{itemize item}{fg=mmgreenDeep}
\setbeamercolor{itemize subitem}{fg=mmgreenDeep}

\newtcolorbox{MMCard}[2][]{enhanced,breakable=false,colback=#2,colframe=black,boxrule=1.5pt,arc=8pt,left=5pt,right=5pt,top=4pt,bottom=4pt,before skip=0pt,after skip=0pt,#1}
\newcommand{\KoruIcon}{\raisebox{-0.2em}{\includegraphics[height=1.05em]{../../../manamaths/SITE/header-logo.png}}}
\newcommand{\NotesTitle}[1]{{\colorbox{mmtanSoft}{\parbox{0.975\linewidth}{\vspace{0.14em}\hspace{0.25em}{\LARGE\bfseries\textcolor{mmgreenDeep}{#1}\hfill\KoruIcon}\vspace{0.14em}}}\par\vspace{0.18em}{\color{mmgreenLeaf}\rule{\linewidth}{2.0pt}}\par\vspace{0.12em}}}
\newcommand{\SectionTitle}[1]{{\large\bfseries\textcolor{mmgreenDeep}{#1}}\par\vspace{0.04em}}
\newcommand{\GreenDivider}{\par\vspace{0.01em}{\color{mmgreenLeaf}\rule{\linewidth}{2.4pt}}\par\vspace{0.03em}}
\newcommand{\KeyIdea}[1]{\begin{MMCard}[title={\textbf{Key idea}},colbacktitle=mmgreenSoft,coltitle=mmgreenDeep]{mmtanSoft}\small #1\end{MMCard}}
\newcommand{\WorkedExample}[2]{\begin{MMCard}[title={\textbf{#1}},colbacktitle=mmgreenSoft,coltitle=mmgreenDeep]{white}\small #2\end{MMCard}}
\newcommand{\CommonMistake}[1]{\begin{MMCard}[title={\textbf{Common mistake}},colbacktitle=mmwarnSoft,coltitle=mmwarn]{white}\small #1\end{MMCard}}
\newcommand{\TryThese}[1]{\begin{MMCard}[title={\textbf{Try these}},colbacktitle=mmgreenSoft,coltitle=mmgreenDeep]{mmtanSoft}\small #1\end{MMCard}}
\newcommand{\StepLine}[2]{\textbf{#1.} #2\par\vspace{0.03em}}

\usepackage{tikz}
\setbeamertemplate{background}{
 \begin{tikzpicture}[remember picture,overlay]
 \draw[black,line width=1.5pt,rounded corners=10pt]
 ([xshift=0.45cm,yshift=-0.45cm]current page.north west)
 rectangle
 ([xshift=-0.45cm,yshift=0.45cm]current page.south east);
 \end{tikzpicture}
}

\begin{document}

\begin{frame}[t,shrink=10]
\NotesTitle{Notes \& Steps}
\footnotesize
\begin{columns}[T,onlytextwidth]
\begin{column}{0.48\textwidth}
\KeyIdea{Weight is measured in grams (g) and kilograms (kg). 1000 g = 1 kg. Everyday items have typical weights you can use as benchmarks.}
\SectionTitle{Steps}
\StepLine{1}{Think about similar everyday items. A paperclip ≈ 1 g. An apple ≈ 200 g. A bag of sugar ≈ 1 kg.}
\StepLine{2}{Decide: is the item lighter or heavier than a bag of sugar (1 kg)?}
\StepLine{3}{Choose the unit: grams for small items (under 1 kg), kilograms for heavy items.}
\StepLine{4}{Check your estimate: does your answer make physical sense?}
\end{column}
\begin{column}{0.48\textwidth}
\SectionTitle{Weight benchmarks}
\begin{itemize}
  \item \textbf{1 g}: a paperclip, a staple, a grain of rice
  \item \textbf{100 g}: a chocolate bar, a small apple
  \item \textbf{200 g}: a cup of flour, a sliced bread loaf
  \item \textbf{500 g}: a block of butter, a small pillow
  \item \textbf{1 kg}: a bag of sugar, a litre of water
  \item \textbf{5 kg}: a school bag full of books
  \item \textbf{10--30 kg}: a large dog, a small child
  \item \textbf{50--100 kg}: an adult, a fridge, a sofa
\end{itemize}
\GreenDivider
\CommonMistake{Thinking everything light uses grams. A TV weighs about 5 kg — not 5000 g or 5 g. Always imagine the item compared to a bag of sugar (1 kg).}
\end{column}
\end{columns}
\end{frame}

\begin{frame}[t,shrink=12]
\NotesTitle{Notes \& Steps}
\begin{columns}[T,onlytextwidth]
\begin{column}{0.48\textwidth}
\WorkedExample{Example 1: a book}{A paperback book: is it measured in grams or kilograms? A book feels lighter than a bag of sugar (1 kg) but heavier than an apple (200 g). It weighs about \textbf{300--500 g}.}
\end{column}
\begin{column}{0.48\textwidth}
\WorkedExample{Example 2: a fridge}{A fridge: grams or kilograms? Much heavier than 1 kg. Compare to an adult (50--80 kg). A fridge weighs about \textbf{50--80 kg}.}
\end{column}
\end{columns}
\GreenDivider
\begin{columns}[T,onlytextwidth]
\begin{column}{0.48\textwidth}
\WorkedExample{Example 3: choosing the right unit}{A feather: 2 grams or 2 kilograms? A feather is lighter than a paperclip (1 g). \textbf{2 g} is correct. 2 kg would be heavier than a bag of sugar — no feather weighs that much.}
\end{column}
\begin{column}{0.48\textwidth}
\WorkedExample{Example 4: comparing weights}{A cat weighs 4 kg. An apple weighs 200 g. How many apples equal one cat? Convert: 4 kg = 4000 g. \(4000 \div 200 = 20\) apples.}
\end{column}
\end{columns}
\end{frame}

\end{document}
